\documentclass[instructions.tex]{subfiles}
\begin{document}
 
\chapter{L'aveuglement de ceux qui négligent leur salut}
 

\primo Tout ce que l'Homme-Dieu a souffert pour notre salut, tout ce que les saints ont fait pour être eux-mêmes sauvés, doit nous apprendre que le salut de notre âme est d'une telle importance, qu'on devrait être prêt à tout sacrifier pour la sauver; mais, par un étrange aveuglement, on dirait qu'on ne travaille que pour la rendre malheureuse. 

L'homme sage ne goûte de véritables plaisirs ici-bas, et n'y trouve rien de solide que ce qui procure le salut; mais l'homme insensé s'aveugle jusqu'à ne trouver de plaisir qu'à se perdre. Il y aurait de la folie d'exposer son âme et de la risquer pour devenir maitre du monde entier, et tous les jours on l'expose à la tentation pour des bagatelles, pour un plaisir frivole, pour plaire à une créature, pour le vil intérêt d'un héritage, ou pour une fumée d'honneur. 

\secundo Si du moins on avait pour son âme autant d'empressement que pour le reste; mais, par un aveuglement profond, on préfère tout à son âme. Fait-on pour la sauver ce qu'on fait pour s'enrichir, et pour conserver la santé du corps ? Que de soins pour sa fortune, pour l'entretien d'une vie misérable qui ne durera que peu de jours, tandis qu'à peine donnera-t-on quelques momens à son âme ! Cette pensée fit une telle impression sur l'esprit d'un secrétaire d'État à la mort, qu'il s'écria en soupirant: \og Oh! que j'ai été insensé, j'ai écrit plus de vingt rames de papier pour le service de mon prince, et je n'ai pas écrit une ligne pour le salut de mon âme ! \fg 

Est-on malade, on s'inquiète, on s'afflige. Est-on en péché mortel, on n'est ni touché, ni affligé; on se divertit, on ne pense ni à son âme, qui est morte, ni à Dieu, qui la doit juger: les semaines, les années, s'écoulent sans qu'on pense à la tirer de l'abime. On choisit le plus habile médecin pour guérir le corps; mais pour l'âme, on choisit souvent un confesseur qui, loin de la guérir, la laisse dans la langueur de la mort. On s'applique à conserver un corps de péché qui doit bientôt pourrir, et à peine pense-t-on à sanctifier une âme immortelle, qui doit régner toujours. 

On veut que tout ce qui sert au corps et à la vie, la nourriture, les meubles, les domestiques, soient excellens; mais pour l'âme, tout paraît indifférent. On dirait qu'elle n'est dans le corps que pour en être l'esclave: elle languit pendant que le corps a ses plaisirs; elle est morte, souillée de crimes, pendant que le corps est flatté. O aveuglement! Un homme ne serait-il pas insensé, qui aurait plus de soin de ses habits que de son corps? Et tous les jours on en voit qui ont plus de soin de leur vêtemens, de leur animaux et de leurs terres, que de leurs âmes. O enfans des hommes! jusqu'à quand aurez-vous le coeur appesanti ? Pourquoi aimer tant la vanité des choses de la terre ?

A quoi aboutissent vos empressemens pour les biens, pour les plaisirs et les honneurs du monde ? Serait-ce donc une grande fortune pour vous d'être en possession de toues les richesses de la terre, si par là vous deviez être damné ? Est-ce aujourd'hui un sujet de consolation à ceux qui sont dans les feux de l'autre vie, d'avoir vécu en celle-ci dans l'abondance et les délices ? Hélas ! quand ils auraient été les maîtres du monde, ils n'en sont à présent que plus malheureux. 

Apprenez à penser sur l'affaire du salut, comme un homme sensé: méditez tous les jours ces paroles de Jésus-Christ: Que sert à l'homme de gagner l'univers s'il perd son âme\footnote{Quid prodest homini, si mundum universum lucretur, animae verò suae detrimentum patiatur. Matth.16.} Peut-on dire qu'on gagne quand on perd son âme ? dit saint Eucher, puisqu'on a tout perdu quand elle est perdue\footnote{Ubi salutis damnum est, illic utiquè jàm lucrum nullum est . S. Euch}.
\section{Résolutions}

\primo Mon salut sera désormais l'affaire dont la considération réglera toutes les autres affaires de ma vie. 

\secundo Et, à l'exemple du grand Apôtre, je regarderai à l'avenir le monde avec ses honneurs, ses biens, et ses plaisirs, comme de la boue.

\prayer{O mon Dieu! détachez-moi des choses qui passent, et remplissez mon coeur de l'amour des biens éternels.}

\end{document}
